%%%%%%%% MTH 4330 EXAMPLE LATEX SUBMISSION FILE %%%%%%%%%%%%%%%%%
%%%%%%%% Adapted from the ICML 2021 EXAMPLE LATEX SUBMISSION FILE %%%%%%%%%%%%%%%%%

\documentclass{article}

% Recommended, but optional, packages for figures and better typesetting:
\usepackage{microtype}
\usepackage{graphicx}
\usepackage{subfigure}
\usepackage{booktabs} % for professional tables

% hyperref makes hyperlinks in the resulting PDF.
% If your build breaks (sometimes temporarily if a hyperlink spans a page)
% please comment out the following usepackage line and replace
% \usepackage{icml2021} with \usepackage[nohyperref]{icml2021} above.
\usepackage{hyperref}

% Attempt to make hyperref and algorithmic work together better:
\newcommand{\theHalgorithm}{\arabic{algorithm}}

% Use the following line for the initial blind version submitted for review:
%\usepackage{icml2021}

% If accepted, instead use the following line for the camera-ready submission:
\usepackage[accepted]{icml2021}

\begin{document}

\twocolumn[
\icmltitle{Multilayer Perceptrons: Mathematical Foundations,\\
a From-Scratch Implementation, and a Three-Dataset Study}

\begin{icmlauthorlist}
\icmlauthor{Leo Chen}{}
\icmlauthor{Sai Kaung San}{}
\icmlauthor{Priyanshu Rana}{}
\end{icmlauthorlist}

\vskip 0.3in
]

% Preliminary report: <= 4 pages, excluding references.
% Section list and required content follow the course guidelines ("Preliminary report").
% An abstract is only needed for the final report.

\section{Introduction}
% - Define the problem class and introduce the selected algorithm.
% - Explain why the topic is suitable for a mathematical and computational study.
% - Note any changes from the approved plan.
% - Put the public GitHub repository link at the end of this section (\url{...}).

\section{Mathematical Foundations}
% - Introduce the notation, objective, assumptions, and central ideas behind the algorithm.
% - Outline the derivation of its main steps.
% - State relevant theoretical properties or limitations.

\section{Implementation Plan and Status}
% - Describe the design of the implementation in Python.
% - Include pseudocode; identify the core components the team will write.
% - Explain how correctness will be checked.
% - State which general-purpose libraries or numerical routines are used.
% - Describe unit tests, synthetic examples, or comparisons with a trusted implementation.
% - Provide the link to the public GitHub repository.

\section{Datasets, Baselines, and Evaluation}
% - For each dataset: source, size, variables, learning task.
% - Preprocessing; data splits or cross-validation; evaluation metrics.
% - Baseline methods; planned tuning procedure.

\section{Progress and Remaining Work}
% - Completed work, pilot results, current difficulties.
% - Schedule and division of responsibilities for the remaining derivations, implementation,
%   experiments, writing, and presentation preparation.
% - Full experimental results are not required at this stage.

% The approved proposal's references are listed with \nocite until the text cites them.
\nocite{rumelhart1986learning,glorot2010understanding,telgarsky2016benefits}

\bibliography{main}
\bibliographystyle{icml2021}

\end{document}
